% !TEX root =./ECE448_Practice_main.tex
%
% L17 practice set (RF reverse engineering). Every question is a step of the
% L17 process on a device of its own: finding the carrier, reading a SigMF
% meta file, turning chips into bits and fields, setting an encoder's blocks,
% and naming a modulation from what the displays show.
% Blank copy: \begin{solution}[h] reserves h of work space.
% SOLUTIONS copy (\iskey=1) prints the red worked solutions.

\fullwidth{\textbf{Learning Objective 6.6: } I can perform RF reverse engineering on
an unknown signal: identify its modulation, symbol rate, and payload structure.
\\[4pt]\normalfont\small Show your work. Write a \textbf{Documentation} line at the
end (\textbf{None} if you did not collaborate).}

\begin{questions}

%%%%%%%%%%%%%%%%%%%%%%%%%%%%%%%%%%%%%%%%%
\question \textbf{Finding the carrier.} A Simulated wireless doorbell is recorded at
1 MS/s, centered at 314.900 MHz, which is an FCC Part 15 frequency. Opened in a viewer
like \texttt{L17-sigmf-view.grc}, the waterfall shows a narrow line 150 kHz below the
center that stays on the whole recording, and short bursts 80 kHz above the center that
appear each time the button is pressed.
\begin{parts}

	\begin{minipage}[t]{0.58\linewidth}
		\part Find the doorbell's carrier frequency, and say how you know which signal is
		the doorbell.
		\begin{solution}[1.0in]
		The bursts appear and stop with each press, so they are the doorbell; the steady
		line belongs to something that is always on.
		\eq{ f_c = 314.900 + 0.080 = 314.980\ \text{MHz} }
		\end{solution}
	\end{minipage}\hfill%
	\begin{minipage}[t]{0.37\linewidth}
		\raggedleft \vspace{12pt}
		\ansbox[1.4in]{314.980 MHz}
	\end{minipage}

	\begin{minipage}[t]{0.58\linewidth}
		\part Give the Frequency Xlating FIR Filter's Center Frequency that shifts the
		doorbell to 0 Hz.
		\begin{solution}[0.7in]
		The filter's center frequency is the offset from the recording's center that it
		moves to 0 Hz: $+80{,}000$ Hz.
		\end{solution}
	\end{minipage}\hfill%
	\begin{minipage}[t]{0.37\linewidth}
		\raggedleft \vspace{12pt}
		\ansbox[1.4in]{80000}
	\end{minipage}

	\part The bursts look several kilohertz wide on the spectrum, and the steady line
	looks almost as thin as one FFT bin. In words, explain why.
		\begin{solution}[0.9in]
		Switching a carrier on and off spreads its power over a band of frequencies around
		it, and the faster it switches, the wider the band. The steady line never changes,
		so its power stays at one frequency.
		\end{solution}

\end{parts}

\newpage
%%%%%%%%%%%%%%%%%%%%%%%%%%%%%%%%%%%%%%%%%
\question \textbf{The meta file.} A teammate sends a SigMF recording. Its
\texttt{.sigmf-meta} file holds these lines:

\begin{center}
\texttt{"core:datatype": "ci16\_le"}, \quad \texttt{"core:sample\_rate": 2000000.0}, \quad
\texttt{"core:frequency": 433000000.0}
\end{center}

\begin{parts}

	\part Give the SigMF Source (Minimal) block's Output Type, the block that must
	follow it, and the flowgraph's \texttt{samp\_rate} and \texttt{fc}.
		\begin{solution}[1.0in]
		The SigMF Source's documentation says \texttt{ci16\_le} is read as short, followed
		by an \textbf{IShort To Complex} block, which pairs the shorts into complex
		samples. \texttt{samp\_rate} = 2000000 and \texttt{fc} = 433e6.
		\end{solution}

	\begin{minipage}[t]{0.58\linewidth}
		\part \texttt{ci16\_le} stores each sample as two 16-bit numbers. How many bytes
		does a 5 s recording take, and how many would \texttt{cf32\_le} take?
		\begin{solution}[1.1in]
		\eq{ 5 \times 2{,}000{,}000 \times 4\ \text{bytes} &= 40\ \text{MB} \\
		5 \times 2{,}000{,}000 \times 8\ \text{bytes} &= 80\ \text{MB} }
		\end{solution}
	\end{minipage}\hfill%
	\begin{minipage}[t]{0.37\linewidth}
		\raggedleft \vspace{12pt}
		\ansbox[1.4in]{40 MB; 80 MB}
	\end{minipage}

	\part URH reads a file it does not recognize as complex float. Explain in words what
	URH would show if you opened this \texttt{ci16\_le} file without converting it.
		\begin{solution}[0.9in]
		URH would read every 8 bytes as one complex float sample, but each 8 bytes holds
		two \texttt{ci16} samples. The values would be meaningless, and the display would
		show noise-like garbage instead of bursts, with half as many samples as the file
		holds. The file has to be converted to \texttt{cf32\_le} first, or opened with an
		extension URH recognizes for 16-bit integers.
		\end{solution}

\end{parts}

\newpage
%%%%%%%%%%%%%%%%%%%%%%%%%%%%%%%%%%%%%%%%%
\question \textbf{Chips to bits.} URH opens the doorbell's recording at 1 MS/s, chooses
ASK and 300 samples per symbol, and lists every burst as the same line of chips:

\begin{center}
\texttt{1000100010101010001010001}
\end{center}

\begin{parts}

	\begin{minipage}[t]{0.58\linewidth}
		\part Find the chip length in microseconds.
		\begin{solution}[0.6in]
		\eq{ \frac{300}{1{,}000{,}000\ \text{samples/s}} = 300\ \mu\text{s} }
		\end{solution}
	\end{minipage}\hfill%
	\begin{minipage}[t]{0.37\linewidth}
		\raggedleft \vspace{12pt}
		\ansbox[1.4in]{300 $\mu$s}
	\end{minipage}

	\part Every pulse is one chip long, and the gaps after the pulses are one chip or
	three chips. Name the encoding, and explain why the line ends with a lone pulse.
		\begin{solution}[1.0in]
		This is \textbf{pulse-distance} coding: every pulse is the same, and the length of
		the gap after it carries the bit. A gap is measured from one pulse to the next, so
		the last bit's gap needs a pulse after it to end; the lone \texttt{1} at the end
		closes the last gap.
		\end{solution}

	\begin{minipage}[t]{0.58\linewidth}
		\part Take a long gap, \texttt{1000}, as a 1 and a short gap, \texttt{10}, as a 0.
		Decode the line into 8 bits and write them in hex.
		\begin{solution}[1.2in]
		Split the line after each pulse's gap:
		\texttt{1000 1000 10 10 10 1000 10 1000 1}, which is 1 1 0 0 0 1 0 1 and a final
		pulse: $1100\ 0101_2 = $ 0xC5.
		\end{solution}
	\end{minipage}\hfill%
	\begin{minipage}[t]{0.37\linewidth}
		\raggedleft \vspace{12pt}
		\ansbox[1.4in]{0xC5}
	\end{minipage}

	\part Find how long this frame lasts, and explain in words why a pulse-distance
	remote has no single bit rate.
		\begin{solution}[1.0in]
		The line is 25 chips: $25 \times 300\ \mu\text{s} = 7.5$ ms. A 1 takes 4 chips and
		a 0 takes 2, so a frame of mostly 1s takes longer than one of mostly 0s; only the
		chip rate, $1/300\ \mu\text{s} \approx 3333$ baud, is fixed.
		\end{solution}

\end{parts}

\newpage
%%%%%%%%%%%%%%%%%%%%%%%%%%%%%%%%%%%%%%%%%
\question \textbf{Encoder settings.} A Simulated garage remote uses pulse-width coding
with 4 chips per bit, \texttt{1110} for a 1 and \texttt{1000} for a 0, with 350
$\mu$s chips. Each frame is 12 bits followed by 10 ms of silence. You build its encoder
as the L17 lesson did, at $f_s = 1$ MS/s: a Vector Source of the bits, a Map, Unpack K
Bits, Char To Float, a Repeat, and a Stream Mux with a Constant Source of 0.
\begin{parts}

	\begin{minipage}[t]{0.58\linewidth}
		\part Give the Map's table and Unpack K Bits' $K$.
		\begin{solution}[1.0in]
		The Map replaces each bit with the number whose binary form is its chip pattern:
		0 with $1000_2 = 8$ and 1 with $1110_2 = 14$. Unpack K Bits with $K = 4$ writes
		those 4 bits out, most significant first.
		\end{solution}
	\end{minipage}\hfill%
	\begin{minipage}[t]{0.37\linewidth}
		\raggedleft \vspace{12pt}
		\ansbox[1.4in]{[8, 14]; $K$ = 4}
	\end{minipage}

	\begin{minipage}[t]{0.58\linewidth}
		\part Give the Repeat's interpolation and the Stream Mux's two lengths.
		\begin{solution}[1.2in]
		\eq{ \text{Repeat} &= 350\ \mu\text{s} \times 10^6 = 350 \\
		\text{frame} &= 12 \times 4 \times 350 = 16{,}800 \\
		\text{gap} &= 0.010 \times 10^6 = 10{,}000 }
		\end{solution}
	\end{minipage}\hfill%
	\begin{minipage}[t]{0.37\linewidth}
		\raggedleft \vspace{12pt}
		\ansbox[1.6in]{350; [16800, 10000]}
	\end{minipage}

	\part Your encoder's output, opened in URH beside the original recording, shows the
	same lines of chips except that every one of yours is 4 chips longer. In words, give
	the most likely mistake.
		\begin{solution}[0.9in]
		Four chips are one bit, so the encoder sends 13 bits instead of 12: the Vector
		Source holds one bit too many, or the Stream Mux's frame length counts 13 bits.
		Comparing the extra chips with the original shows which bit was added.
		\end{solution}

\end{parts}

\newpage
%%%%%%%%%%%%%%%%%%%%%%%%%%%%%%%%%%%%%%%%%
\question \textbf{Which modulation?} Four unknown recordings open in a viewer like
\texttt{L17-sigmf-view.grc}. Name the modulation of each from what the displays show,
and say which lesson's signal it resembles.
\begin{parts}

	\part The waterfall shows two parallel tracks a few kilohertz apart, and the signal
	is always on one of them; the envelope is constant.
		\begin{solution}[0.7in]
		Frequency-shift keying: the carrier moves between two frequencies, which is
		Lesson 8's FSK.
		\end{solution}

	\part The waterfall shows one band that stays on and does not move, and the envelope
	is nearly constant; a constellation plot after carrier recovery shows four clusters.
		\begin{solution}[0.7in]
		Phase-shift keying, here QPSK: the amplitude and frequency stay put and the phase
		carries the data, which is Lesson 15's mPSK.
		\end{solution}

	\part Short bursts appear at one frequency, and the envelope switches between zero
	and one level in pulses of two lengths.
		\begin{solution}[0.7in]
		On-off keying with pulse-width coding, which is the L17 fan remote, and the
		keyed carrier of Lessons 11 and 12.
		\end{solution}

	\part Short dashes of the same width appear at a new frequency every few hundred
	milliseconds, stepping around a band.
		\begin{solution}[0.7in]
		Frequency hopping, as in Lesson 16: one narrowband signal moving from channel to
		channel. A receiver tuned to one channel would see only the hops that land on its
		channel, so the viewer has to show the whole band.
		\end{solution}

\end{parts}

\newpage
%%%%%%%%%%%%%%%%%%%%%%%%%%%%%%%%%%%%%%%%%
\question \textbf{Repeats and noise.} The L17 fan remote sends each frame 6 times per
press and ends each frame with a check nibble.
\begin{parts}

	\part In words, explain what the repeats and the check each let a receiver do when
	the signal is weak.
		\begin{solution}[1.0in]
		Noise can flip a chip and spoil a frame. The check lets the receiver recognize a
		spoiled frame and ignore it instead of acting on the wrong command, and the repeats
		give it five more chances to receive a frame that passes the check.
		\end{solution}

	\part You lower the signal in the pulse decoder explorer by raising the noise. Name
	two things you would watch to see the frames start to fail, in the explorer or in
	URH.
		\begin{solution}[0.9in]
		In the explorer, groups marked \texttt{?} that are neither \texttt{110} nor
		\texttt{100}, and now and then the Frame readout reporting that the check fails. In URH, the lines
		of chips for one press no longer matching each other from repeat to repeat.
		\end{solution}

\end{parts}

\end{questions}
